\section{Component gerbes and terminal rigidity}
\label{sec:gerbe-obstruction}

For a stack in groupoids \(\mathfrak E\), let \(\pi_0(\mathfrak E)\) denote
the sheafification of its presheaf of fibrewise isomorphism classes.  If
\(v\in\pi_0(\mathfrak E)(a)\), define the full substack
\begin{equation}\label{eq:component-substack}
 \mathfrak E[v](X)=
 \{x\in\mathfrak E(X):[x]=v|_X\}.
\end{equation}
The equality in \eqref{eq:component-substack} is equality in the component
sheaf, so it is local by construction.

The exact ancestry of this construction is classical.  Giraud proves that
the projection of a stack to its sheaf of connected components is a gerbe
over that sheaf, and that pullback along a component section gives the
corresponding maximal subgerbe
\cite[III, Prop.~2.1.5.3]{Giraud1971}.  The next proposition records this
standard construction and the two compatibilities used later.  It is an
input, not an originality claim.

\begin{proposition}[Standard component-gerbe package]
\label{prop:standard-component-gerbe-package}
Let \(\mathfrak E\) be a stack in groupoids on \(\mathcal C\), and let
\(v\in\pi_0(\mathfrak E)(a)\).
\begin{enumerate}[label=(\roman*),nosep]
\item The full substack \(\mathfrak E[v]\) is the maximal gerbe over the
      section \(v\).  If \(\mathfrak E\) is banded by an abelian sheaf
      \(\mathscr A\), then \(\mathfrak E[v]\) has the restricted
      \(\mathscr A\)-banding.
\item Its Giraud class satisfies
      \[
      [\mathfrak E[v]]=0\text{ in }H^2(\mathcal C,\mathscr A)
      \quad\Longleftrightarrow\quad
      \mathfrak E[v]\text{ is neutral}
      \quad\Longleftrightarrow\quad
      \mathfrak E[v](a)\ne\varnothing.
      \]
\item Suppose \(\Phi:\mathfrak E\to\mathfrak E'\) is a banded equivalence,
      its band map is \(\theta:\mathscr A\xrightarrow{\sim}\mathscr A'\),
      and its component map carries \(v\) to \(v'\).  Then \(\Phi\)
      restricts to a banded equivalence
      \(\mathfrak E[v]\simeq\mathfrak E'[v']\), and
      \[
      \theta_*[\mathfrak E[v]]=[\mathfrak E'[v']].
      \]
\end{enumerate}
\end{proposition}

\begin{proof}
Clause (i) is precisely the component-gerbe construction of
\cite[III, Prop.~2.1.5.3]{Giraud1971}.  Concretely, \(v|_X\) is locally
represented by an object, so \(\mathfrak E[v]\) is locally nonempty; two of
its objects have the same component class and are therefore locally
isomorphic.  Fullness preserves the automorphism sheaves and hence the band.

The classification of abelian-banded gerbes identifies equivalence classes
with \(H^2\), with the neutral gerbes representing zero
\cite[Chap.~IV]{Giraud1971}.  On a site with terminal object \(a\), neutrality
is exactly existence of an object over \(a\), proving (ii).  For (iii), the
component compatibility makes \(\Phi\) restrict to the full component
substacks.  Applying \(\Phi\) to local objects and overlap arrows applies
\(\theta\) to their triple-overlap defect cocycle.  This is the standard
naturality of Giraud classification under a banded equivalence
\cite[Tag 06NY]{StacksProject}.
\end{proof}

Apply this package to the stackification
\(\mathfrak L=a(\mathfrak P)\) of \Cref{sec:preliminaries}.  In view of
\eqref{eq:stackification-components-core}, every
\(v\in(aF)(a)\) determines the component gerbe
\[
 \mathfrak G_v:=\mathfrak L[v]
\]
and, when \(\mathfrak P\) is \(\mathscr A\)-banded, a class
\[
 \omega_v:=[\mathfrak G_v]\in H^2(\mathcal C,\mathscr A).
\]
On a Leray cover, choosing local objects and overlap arrows produces the
usual Cech \(2\)-cocycle; changing either choice changes it by a coboundary.
This is the local representative of the same Giraud class, not a second
obstruction construction \cite{Breen2006,Moerdijk2002}.

\subsection*{The standard cocycle in local coordinates}

We spell out the preceding sentence because the realization theorem will use
the same coordinates in the opposite direction.  Let
\(\mathcal U=\{U_i\to a\}\) be a cover on which
\(\mathfrak G_v\) has objects \(x_i\in\mathfrak G_v(U_i)\).  After refining
the cover when necessary, choose isomorphisms
\[
 g_{ij}:x_i|_{U_{ij}}\longrightarrow x_j|_{U_{ij}},
 \qquad
 g_{ii}=1,\quad g_{ji}=g_{ij}^{-1}.
\]
On a triple overlap \(U_{ijk}\), the two arrows from \(x_i\) to \(x_k\)
need not agree.  Their ratio is an automorphism of \(x_i\):
\begin{equation}\label{eq:giraud-triple-defect}
 c_{ijk}
 :=g_{ik}^{-1}g_{jk}g_{ij}
 \in\underline{\operatorname{Aut}}(x_i)(U_{ijk})
 \cong\mathscr A(U_{ijk}).
\end{equation}
The last identification is the chosen banding.  We write the band
additively.  Comparing the five composites around a quadruple overlap gives
\begin{equation}\label{eq:giraud-cocycle-equation}
 c_{jkl}-c_{ikl}+c_{ijl}-c_{ijk}=0.
\end{equation}
Thus \(c=(c_{ijk})\) is a Cech \(2\)-cocycle.

There are two relevant changes of choice.  First, keep the local objects
fixed and replace the overlap arrows by
\[
 g'_{ij}=b_{ij}\cdot g_{ij},
 \qquad b_{ij}\in\mathscr A(U_{ij}).
\]
Because the band is abelian, substituting in
\eqref{eq:giraud-triple-defect} gives
\[
 c'_{ijk}
 =c_{ijk}+b_{jk}-b_{ik}+b_{ij}
 =c_{ijk}+(\delta b)_{ijk}.
\]
Second, replace \(x_i\) by an isomorphic local object \(x'_i\) and choose
arrows \(h_i:x_i\to x'_i\).  Transporting \(g_{ij}\) by the \(h_i\) gives
overlap arrows for the new objects with the same defect under the band
identification.  Any other choice of transported overlap arrows differs
from these by a \(1\)-cochain and therefore changes the defect by a
coboundary.  Refining the cover pulls the cocycle back.  The resulting
class is consequently independent of all local choices in the Cech colimit,
and the Leray comparison sends it to
\([\mathfrak G_v]\in H^2(\mathcal C,\mathscr A)\).

The neutral case is visible in the same coordinates.  A global object
\(x\in\mathfrak G_v(a)\) restricts to local objects whose canonical overlap
arrows have zero triple defect.  Conversely, if the defect of a chosen local
system is a coboundary, modifying its overlap arrows by the corresponding
\(1\)-cochain makes the defect zero.  Effective descent in the stack then
glues the local objects to a global object.  This is the cocycle-level
description of the neutrality equivalence in
\Cref{prop:standard-component-gerbe-package}(ii); it is not a replacement
for Giraud classification on a general site.

Finally, if a banded equivalence has band map
\(\theta:\mathscr A\to\mathscr A'\), it sends the chosen local objects and
arrows to local objects and arrows whose defect is \(\theta(c_{ijk})\).
This makes the naturality in
\Cref{prop:standard-component-gerbe-package}(iii) literal on a common
presentation.

\begin{theorem}[Presheaf-to-component-gerbe comparison criterion]
\label{thm:gerbe-null-semantics}
Let \(F\) be separated on the slice over \(a\), and let
\(\mathfrak P\) be an abelian-banded prestack lift whose stackification map
is terminally essentially surjective.  Then
\begin{equation}\label{eq:matching-gerbe-criterion}
 \operatorname{MWA}(F,a)
 \quad\Longleftrightarrow\quad
 (aF)(a)\ne\varnothing
 \ \text{and}\
 \mathfrak G_v\text{ is non-neutral for every }v\in(aF)(a).
\end{equation}
Equivalently, every \(\omega_v\) is nonzero.  The dependencies are
\[
\begin{array}{c@{\quad\Longleftrightarrow\quad}c}
 \operatorname{Match}(F,a)&(aF)(a)\ne\varnothing
 \quad\text{(slice separatedness)},\\[2mm]
 v\in(aF)(a)&\mathfrak G_v\text{ a component gerbe}
 \quad\text{(Giraud III.2.1.5.3)},\\[2mm]
 F(a)\ne\varnothing&\mathfrak L(a)\ne\varnothing
 \quad\text{(terminal essential surjectivity)},\\[2mm]
 \mathfrak G_v(a)\ne\varnothing&\mathfrak G_v\text{ neutral}
 \quad\text{(Giraud classification)}.
\end{array}
\]
\end{theorem}

\begin{proof}
By \Cref{prop:finite-site-semantic-interface}(i), matching is equivalent to
\((aF)(a)\ne\varnothing\).  By clause (ii) of that proposition and terminal
essential surjectivity, \(F(a)\ne\varnothing\) is equivalent to
\(\mathfrak L(a)\ne\varnothing\).  The latter holds exactly when one of the
full components \(\mathfrak G_v\) contains an object over \(a\).  By
\Cref{prop:standard-component-gerbe-package}(ii), this is equivalent to
neutrality of at least one component gerbe.  Negating that statement while
retaining nonemptiness of \((aF)(a)\) proves
\eqref{eq:matching-gerbe-criterion}; the cohomological form follows from the
same standard package.
\end{proof}

The criterion is a site-theoretic reformulation of the meet-semilattice
comparison in the companion manuscript
\cite[Theorem~5.20]{MaZhang2026ConditionalGluingFailure}; it is not an
independent gerbe-classification theorem.  Its role here is to connect the
strict presheaf predicate to the standard component-gerbe construction on a
pullback site.

We now give the representative-rigidity argument in full.  This result
isolates a restriction on attempts to realize matching without amalgamation
by an excessively strict lift.

\begin{theorem}[Terminal rigidity for representative-rigid lifts]
\label{thm:canonical-lift-rigidity}
Let \(F\) be separated on the slice over \(a\), and let \(\mathfrak P\) be a
representative-rigid \(\mathscr A\)-banded prestack lift of \(F\).  Write
\(u:F\to aF\) and \(\iota:\mathfrak P\to\mathfrak L=a(\mathfrak P)\).
Then:
\begin{enumerate}[label=(\roman*),nosep]
\item terminal essential surjectivity of \(\iota\) implies that
      \(u_a:F(a)\to(aF)(a)\) is surjective;
\item if \(H^1(\mathcal C,\mathscr A)=0\), surjectivity of \(u_a\) implies
      terminal essential surjectivity of \(\iota\).
\end{enumerate}
Under the same \(H^1\)-vanishing hypothesis, the following are equivalent:
\begin{enumerate}[label=(\alph*),nosep]
\item \(\iota\) is terminally essentially surjective;
\item \(u_a\) is bijective;
\item every strict matching family on every cover of \(a\) has a unique
      amalgamation in \(F(a)\).
\end{enumerate}
Strict pullback stability of the selected representatives is used in
\textup{(i)}.  The converse \textup{(ii)} uses only the standard
\(\mathscr A\)-torsor of isomorphisms and its \(H^1\)-class.
\end{theorem}

\begin{proof}
We separate the proof into the forward descent construction, the converse
torsor argument, and the comparison with strict amalgamation.

\medskip
\noindent\emph{Forward descent construction.}
Let
\(\xi\in(aF)(a)=\pi_0(\mathfrak L)(a)\).  By the proof of
\Cref{prop:finite-site-semantic-interface}(i), there is a cover
\(\{U_i\to a\}\) and a strict matching family \(t_i\in F(U_i)\)
representing \(\xi\).  Representative-rigidity supplies objects
\[
 y_i=s_{U_i}(t_i)\in\mathfrak P(U_i).
\]
On \(U_{ij}\), strict pullback stability gives literal equalities
\[
 y_i|_{U_{ij}}
 =s_{U_{ij}}(t_i|_{U_{ij}})
 =s_{U_{ij}}(t_j|_{U_{ij}})
 =y_j|_{U_{ij}}.
\]
The identity arrows of these common objects are therefore legitimate
descent arrows.  Their triple composites are identities on the nose, so
stack descent glues the \(y_i\) to an object
\(x_\xi\in\mathfrak L(a)\) with component class \(\xi\).  This is the sole
place where strict pullback stability is used.  With merely
pseudofunctorial representatives, cartesian comparison arrows would enter
and their coherence would have to be controlled separately.

The descent datum just constructed is stronger than the datum provided by
equality of component classes alone.  Equality
\([y_i|_{U_{ij}}]=[y_j|_{U_{ij}}]\) would only guarantee local
isomorphisms after further covers.  Equation
\eqref{eq:strict-representative-pullback} instead gives the same object on
the overlap.  Hence the chosen descent arrow is the identity, not a selected
arrow in a nontrivial Isom torsor.  On \(U_{ijk}\) the descent equation is
\[
 1_{y_k}\circ1_{y_j}=1_{y_i}
\]
after the literal identifications of the three restrictions.  No
\(\mathscr A\)-valued \(2\)-cocycle remains to be killed.  This explains
exactly why strict pullback stability, rather than a componentwise choice of
objects, is the forward hypothesis.

Assume terminal essential surjectivity.  There is
\(y\in\mathfrak P(a)\) and an isomorphism
\(\iota(y)\cong x_\xi\).  Let \(t\in F(a)\) be the component of \(y\).
The component map of \(\iota\) is \(u\), hence
\[
 u_a(t)=[\iota(y)]=[x_\xi]=\xi.
\]
Since \(\xi\) was arbitrary, \(u_a\) is surjective.  This proves (i).

\medskip
\noindent\emph{The converse Isom-torsor argument.}
For (ii), assume \(u_a\) is surjective and take any
\(x\in\mathfrak L(a)\).  Choose \(t\in F(a)\) with
\(u_a(t)=[x]\), and put \(y=s_a(t)\).  The sheaf
\[
 \mathcal I_{y,x}:=
 \underline{\operatorname{Isom}}_{\mathfrak L}(\iota(y),x)
\]
is locally nonempty because the two objects have the same component class.
Composition makes it a simply transitive torsor under
\(\underline{\operatorname{Aut}}(\iota(y))\cong\mathscr A\).  This is the
usual Isom-torsor of a banded stack
\cite[Chap.~III]{MacLaneMoerdijk1994}, \cite[Tag 04TU]{StacksProject}.
More explicitly, choose a cover \(\{V_j\to a\}\) and isomorphisms
\[
 h_j:\iota(y)|_{V_j}\xrightarrow{\sim}x|_{V_j}.
\]
On \(V_{jk}\), the ratio
\[
 d_{jk}=h_k^{-1}h_j
\]
is an automorphism of \(\iota(y)|_{V_{jk}}\), hence a section of
\(\mathscr A(V_{jk})\).  The equality of threefold composites gives
\[
 d_{jk}+d_{kl}=d_{jl}
\]
on \(V_{jkl}\), so \(d=(d_{jk})\) is a Cech \(1\)-cocycle representing
the Isom torsor.  If its class vanishes, there are sections
\(e_j\in\mathscr A(V_j)\) with \(d_{jk}=e_k-e_j\).  Replacing
\(h_j\) by \(h_j\cdot e_j\) makes the local isomorphisms agree on overlaps,
and the sheaf condition for Isom glues them to a global isomorphism
\(\iota(y)\cong x\).

The hypothesis \(H^1(\mathcal C,\mathscr A)=0\) makes every such torsor
trivial, so \(\mathcal I_{y,x}\) has a global section.  Thus
\(\iota(y)\cong x\), proving terminal essential surjectivity.

\medskip
\noindent\emph{Comparison with strict amalgamation.}
It remains to compare (b) and (c).  Slice separatedness implies that \(u_a\)
is injective: equality after sheafification is witnessed after a cover, and
separatedness at \(a\) returns equality in \(F(a)\).  Thus surjectivity and
bijectivity of \(u_a\) are equivalent.

Suppose \(u_a\) is bijective and let \((t_i)\) be a strict matching family
on a cover of \(a\).  It glues in \(aF\) to \(\xi\); choose
\(t\in F(a)\) with \(u_a(t)=\xi\).  Over each \(U_i\), the sections
\(t|_{U_i}\) and \(t_i\) become equal in \(aF(U_i)\).  Such equality is
witnessed after a further cover of \(U_i\), and separatedness at that object
of the slice gives \(t|_{U_i}=t_i\).  Hence \(t\) is an amalgamation, unique
by separatedness at \(a\).

Conversely, assume every matching family on a cover of \(a\) has a unique
amalgamation.  Given \(\xi\in(aF)(a)\), the plus--plus representation and
slice separatedness used in
\Cref{prop:finite-site-semantic-interface}(i) produce a strict matching
family representing \(\xi\).  Its amalgamation \(t\in F(a)\) satisfies
\(u_a(t)=\xi\), so \(u_a\) is surjective and hence bijective.  Combining
this equivalence with (i) and (ii) proves (a)--(c).

The roles of the hypotheses are now separate.  Slice separatedness upgrades
locally witnessed equality of presheaf sections to strict equality and gives
uniqueness of amalgamations.  Representative-rigidity turns a strict
matching family into identity-valued object descent data.  Terminal
essential surjectivity is used only in the forward passage from a stack
object to a prestack object.  Finally,
\(H^1(\mathcal C,\mathscr A)=0\) is used only to trivialize the Isom torsor
in the converse.  None of these steps supplies a new classification of
gerbes.
\end{proof}

Consequently, on a slice with \(H^1(\mathcal C,\mathscr A)=0\), matching
without amalgamation cannot coexist with both terminal essential
surjectivity and a strictly pullback-stable representative choice.  This is
a no-go statement for that rigid presentation of a lift; it is not a new
general descent or gerbe-classification principle.
